\subsection*{电荷分布的三种连续模型}

当大量微小带电粒子密集分布时，逐个求和既不方便也没有必要。按电荷所占据的几何维数，常用
线电荷密度、面电荷密度和体电荷密度：
\begin{equation}
 \lambda=\frac{\dd q}{\dd l},
 \qquad
 \sigma=\frac{\dd q}{\dd S},
 \qquad
 \rho=\frac{\dd q}{\dd V}.
\end{equation}
总电荷分别由 $Q=\int\lambda\,\dd l$、$Q=\int\sigma\,\dd S$、
$Q=\int\rho\,\dd V$ 得到。选用哪一种模型，取决于带电体厚度与观察尺度的比较；同一根有有限
半径的带电细棒，在远处可近似为线电荷，在极远处又可进一步近似为点电荷。

对连续分布，场点记为 $\vect r$，源点记为 $\vect r'$，电场的统一表达式为
\begin{equation}
 \vect E(\vect r)=\frac1{4\pi\varepsilon_0}
 \int\frac{\vect r-\vect r'}{|\vect r-\vect r'|^3}\,\dd q(\vect r').
\end{equation}
必须区分场点坐标和源点坐标；分母是二者距离的三次方，分子提供从源点指向场点的方向。对称性
可以在积分前消去部分分量。例如场点位于均匀带电圆环轴线上时，所有横向分量成对抵消，只需
积分轴向分量。

叠加原理来自静电场方程的线性。若电荷分布写成几部分之和，电场就是各部分电场的矢量和。
“挖去空腔”常用负电荷叠加：先填满整个带电体，再在空腔位置叠加一个形状相同、密度相反的
带电体。这样既保留原问题的几何条件，又能直接使用已知的高对称场。

\subsection*{电场线、试探电荷与粒子轨迹}

试探电荷应足够小，使它对原有电荷分布的扰动可以忽略。定义式 $\vect E=\vect F/q$ 中的 $q$
带有符号，因此无论试探电荷取正取负，所得 $\vect E$ 都相同。教学图中通常想象用正试探电荷
确定电场方向；负电荷受到的力与电场方向相反。

电场线的疏密没有绝对数值含义，只有在同一幅图采用统一画法时，才能定性比较场强。给定一小
面积元，穿过它的电场线数可被安排为与电通量成正比；因此电场线不会在无电荷区域凭空开始或
终止。静电场线也不能形成闭合曲线，因为若沿闭合电场线前进，$\vect E\cdot\dd\vect l$ 始终
同号，会与静电场环流为零矛盾。

带电粒子的轨迹由 Newton 第二定律决定：
\begin{equation}
 m\ddot{\vect r}=q\vect E(\vect r).
\end{equation}
轨迹方向由粒子的初速度和电荷符号共同决定，一般不会与电场线重合。只有正电荷从静止释放，且
运动过程中电场方向不发生横向改变等特殊情况下，轨迹才可能沿电场线。

\subsection*{高斯定律怎样真正用于求场}

高斯定律给出的是闭合面上法向电场的积分，单靠它不能确定曲面每一点的电场。要把积分化成
$ES$，必须同时证明：电场在所选曲面的有效部分大小相等，并且方向与面积元平行或垂直。常见的
三类选择是：
\begin{itemize}
  \item 球对称分布选同心球面，$\vect E$ 沿径向且同一球面上大小相等；
  \item 无限长柱对称分布选同轴圆柱面，侧面贡献 $E(2\pi rl)$，两个端面通量为零；
  \item 无限大平面对称分布选跨过平面的薄柱，两个端面贡献相等，侧面通量为零。
\end{itemize}

以半径 $R$、体电荷密度均匀为 $\rho$ 的实心球为例。球内半径 $r<R$ 的高斯面包围电荷
$Q_r=4\pi r^3\rho/3$，故
\begin{equation}
 E(r)=\frac{\rho r}{3\varepsilon_0}.
\end{equation}
球外则只看总电荷 $Q=4\pi R^3\rho/3$，有
\begin{equation}
 E(r)=\frac{\rho R^3}{3\varepsilon_0r^2}.
\end{equation}
内外结果在 $r=R$ 连续。若表面另有面电荷，法向电场会在表面发生跃变；把跨越表面的薄柱作为
高斯面可得
\begin{equation}
 E_n^{(+)}-E_n^{(-)}=\frac{\sigma}{\varepsilon_0}.
\end{equation}

静电条件 $\nabla\times\vect E=0$ 还给出切向分量连续：
$E_t^{(+)}=E_t^{(-)}$。因此一张理想面电荷会使法向分量跳变，却不会使切向分量跳变。这两个
边界关系在导体、电介质和分区电场问题中会反复使用。

\subsection*{积分形式与微分形式的互补}

积分形式关注一个有限闭合面包围的总电荷，适合利用整体对称性；微分形式关注一点附近的电场
散度，适合研究局部变化。二者完全等价，但并不意味着“某点散度为零就有电场为零”。在没有
电荷的区域，$\nabla\cdot\vect E=0$，点电荷产生的非零电场仍可以穿过该区域，只是任何不包围
电荷的小闭合面的净通量为零。

静电场还满足
\begin{equation}
 \nabla\times\vect E=\vect0,
 \qquad
 \oint\vect E\cdot\dd\vect l=0.
\end{equation}
高斯定律决定电场的“源”，零旋度条件决定它没有静电环流。两条性质和适当边界条件合在一起，
才能唯一确定静电场。
